% TOP_PROBLEM: {"id": 3700043, "title": "Few inflection points of holomorphic curves", "category_id": 8, "category": {"id": 8, "name": "analysis", "display_name": "Analysis", "description": "Limits, continuity, calculus, and function theory.", "slug": "analysis", "order_index": 8, "created_at": "2026-07-31T15:26:25.671Z"}, "status": "partially_solved"}
% FIRST_POSTED: null
% ENTRY_KIND: statement_only
% Imported from: batch12.tex; UnsolvedMath ID 3700043
% Compile twice: lualatex 3700043.tex
\documentclass[11pt]{article}
\usepackage{fontspec}
\setmainfont{Latin Modern Roman}
\usepackage{amsmath,amssymb,mathtools,unicode-math}
\setmathfont{Latin Modern Math}
\usepackage{xurl,hyperref}
\hypersetup{hidelinks,pdftitle={UnsolvedMath Problem 3700043}}
\newcommand{\sref}[2]{\hyperlink{source:#1:#2}{[S#2]}}
\newcommand{\cataloguescope}[1]{\paragraph{Mathematical question.}#1}
\begin{document}
% UnsolvedMath ID 3700043; source catalogue code AMR-036-0043
\setcounter{section}{3700042}
\section{Few inflection points of holomorphic curves}\label{3700043}
{\small\sffamily UnsolvedMath ID 3700043 · Analysis}
\subsection{Definitions and mathematical statement}

\paragraph{Shared notation.}$\mathbb N=\{1,2,\ldots\}$, and $\mathbb Z,\mathbb Q,\mathbb R,\mathbb C$ denote the integers, rationals, reals and complex numbers. Any other conventions are specified locally.


Let $n\ge1$ and let $f:\mathbb C\to\mathbb P^n$ have a reduced entire representation $f=[f_0:\cdots:f_n]$, where $f_0,\ldots,f_n$ have no common zero and are linearly independent over $\mathbb C$. Define
\[
 T(r,f)=\frac{1}{2\pi}\int_0^{2\pi}\log\Bigl(\sum_{j=0}^{n}|f_j(re^{i\theta})|^2\Bigr)^{1/2}\,d\theta
 -\log\Bigl(\sum_{j=0}^{n}|f_j(0)|^2\Bigr)^{1/2}.
\]
Let $W_f=\det(f_j^{(i)})_{0\le i,j\le n}$ be its Wronskian. Write $m_0$ for the zero multiplicity of $W_f$ at $0$, and let $n_1(t)$ count its zeros in $|z|\le t$, with multiplicity. For $r\ge1$, put
\[
 N_1(r)=m_0\log r+\int_0^r\frac{n_1(t)-m_0}{t}\,dt,
 \qquad\lambda=\liminf_{r\to\infty}\frac{\log T(r,f)}{\log r}.
\]
The assumption of linear independence makes $W_f\not\equiv0$. Changing the reduced representation multiplies it by a nonvanishing factor, so the zero count is intrinsic; $T$ changes by at most a harmless normalization constant.
\paragraph{Statement recorded in the supplied source.}
Assume $\lambda<\infty$ and $N_1(r)=o(T(r,f))$ as $r\to\infty$. Must $\lambda$ be rational, and must the limit
\[
 \lim_{r\to\infty}\frac{\log T(r,f)}{\log r}
\]
exist (and thus equal $\lambda$)? The original survey §4 proposes the stronger regularly varying form $T(r,f)=r^\lambda\ell(r)$, with $\ell(cr)/\ell(r)\to1$ uniformly for $1\le c\le2$. The registered question asks for rationality and existence of the growth exponent. \sref{3700043}{2}
\subsection{Short English statement}
For a linearly nondegenerate entire projective curve of finite lower order with negligible Wronskian zero count, prove that its growth exponent exists and is rational.
\subsection{Sources}
\begin{description}
\item[\hypertarget{source:3700043:1}{S1}] ULAM.AI and the UnsolvedMath Contributors, UnsolvedMath: A Curated Collection of Open Mathematics Problems, record 3700043 (AMR-036-0043). CC BY 4.0. \url{https://huggingface.co/datasets/ulamai/UnsolvedMath}.
\item[\hypertarget{source:3700043:2}{S2}] Alexandre Eremenko, Value Distribution and Potential Theory (2003), §§3–4, pp.685–688; §4 Conjecture. \url{https://arxiv.org/abs/math/0304332}.
\end{description}
\label{end:3700043}
\end{document}
